%%%%%%%%%%%%%%%%%%%%%%%%%%%%%%%%%%%%%%%%%%%%%%%%%%%%%%%%%%%%%%%%%%%%%%%%%%%
% Lehreinheit 5: Hypothesen und Falsifizierbarkeit
% Quelle: WuZ-Buch, Kap. 15 (LE 5), Kap. 5b (Popper, Wason/Congruence
%         Bias, Horwich), Kap. 10 (2-4-6-Test), Kap. 10b (Testbarkeit),
%         Anhang E (Kommentar zu Exposure und Grenzarbeit: Switchboard)
% Staunmoment: Switchboard / Black Box  |  Kernbegriff: Falsifizierbarkeit
% Geheimhaltung: Methode des Switchboards wird NICHT beschrieben.
%%%%%%%%%%%%%%%%%%%%%%%%%%%%%%%%%%%%%%%%%%%%%%%%%%%%%%%%%%%%%%%%%%%%%%%%%%%
\nsl{In LE~2 wurde aus einem Staunen eine Forschungsfrage. Jetzt wird aus
der Frage eine Behauptung, die an der Wirklichkeit scheitern kann. Die
Einheit beginnt mit einem Gerät, das jede Ingenieurstudentin nach zehn
Minuten verstanden zu haben glaubt: vier Schalter, vier Lampen, sichtbare
Leiterbahnen. Die Gruppe bildet ein Modell, sammelt Bestätigungen und
erlebt dann, dass beliebig viele Bestätigungen ein falsches Modell nicht
retten. Daraus entwickelt die Einheit die Unterscheidung von Induktion,
Deduktion und Abduktion, Poppers Falsifikationsprinzip, die Formulierung
von Null- und Alternativhypothese, Horwichs graduellen Bestätigungsbegriff
und die Grenzen der Falsifikation (Duhem-Quine). Der Wason-2-4-6-Test
zeigt anschließend, dass der Hang zum Bestätigen eine feste Gewohnheit
ist, gegen die man sich methodisch schützen muss.
\newpage}
\os{\begin{description}
\item[{\rot\bf Heute}]~\\[-\lblskip]
\bi
  \ite ein Schaltbrett, das sich nicht an seine Leiterbahnen hält
  \ite Induktion, Deduktion, Abduktion
  \ite Popper: Falsifikation statt Verifikation
  \ite Null- und Alternativhypothese
  \ite Horwich: Bestätigung hat Grade -- Duhem-Quine: Widerlegung auch
  \ite der 2-4-6-Test und ein Raster für gute Hypothesen
\ei
\end{description}
\newpage}

\begin{vorbereitung}[Switchboard / Black Box]
\textbf{Eine Woche vorher}
\begin{checkliste}
  \item Das \emph{Switchboard} ist ein kommerzielles Kunststück
    \parencite{andersen2020switchboard}: eine offene Platine mit vier
    Schaltern mit farbigen Hülsen und vier Lampen in denselben Farben,
    sichtbar über Leiterbahnen verbunden. Gerät rechtzeitig beim
    Fachhandel für Zauberbedarf beschaffen bzw.\ aus dem Bestand holen.
  \item \textbf{Methode und Einstudierung nach der Anleitung des
    Herstellers; hier bewusst nicht beschrieben.} Das Skript wird
    öffentlich gespiegelt. Assistenzen, die das Gerät mit bedienen,
    werden persönlich eingewiesen und zur Verschwiegenheit verpflichtet.
  \item Batterien prüfen und einen vollständigen Satz Ersatzbatterien
    einpacken; alle Leuchtmittel einzeln testen, Ersatzleuchtmittel
    bereitlegen, falls der Hersteller sie liefert.
  \item Den Ablauf mindestens zweimal vollständig proben, einmal davon
    vor einer Testperson: Exploration durch die Gruppe, eigene
    Behauptung \enquote{Die Farbe der Hülsen ist entscheidend}, Umstecken
    der Hülsen, erneute Exploration, zweiter und dritter Modellwechsel.
    Die Übergänge zwischen den Phasen so einüben, dass sie ohne Blick auf
    Notizen sitzen.
  \item Sprechtext für den Abschluss festlegen (siehe \enquote{Am Tag
    selbst}). Die Haltung dazu steht in Markus Pfeils Kommentar zu
    Exposure und Grenzarbeit (\buch{Anhang~E}): Methode nicht verraten,
    aber die \emph{Richtung} andeuten.
  \item Für den 2-4-6-Test genügt Tafel oder Flipchart. Die Regel
    (\enquote{drei aufsteigende Zahlen}) nur im Kopf haben, nicht auf
    einem Zettel am Pult.
\end{checkliste}

\textbf{Handout und Übungsblatt}
\begin{checkliste}
  \item Handout \enquote{Hypothesen und Falsifizierbarkeit} sowie
    Übungsblatt LE~5 Teil~A und Teil~B je einmal pro Person kopieren
    (Druckvorlage \texttt{handouts.pdf}, Seiten zu LE~5). Teil~A enthält
    die fünf Aussagen für die Klassifikationsübung und liegt zu Beginn
    aus. Teil~B (Rückblick auf das Schaltbrett,
    H\textsubscript{0}/H\textsubscript{1}-Vorlage) setzt mehrere
    Modellwechsel voraus und wird erst in der Anwendungsphase
    ausgeteilt.
  \item Pro Kleingruppe (4--5 Personen) ein A4-Blatt
    \enquote{Protokoll Schaltbrett} mit zwei Spalten \emph{Eingabe} und
    \emph{Beobachtung}; die Gruppen protokollieren ihre Versuche dort.
  \item Kleine Zettel (ca.\ A7) für die stille Regel-Notiz im
    2-4-6-Test, einer pro Person.
  \item Die Forschungsfragen aus LE~2 bereithalten (Rücklauf aus der
    Hausaufgabe); jede Person braucht ihre eigene Frage für Aufgabe~3
    (Übungsblatt Teil~B).
\end{checkliste}

\textbf{Am Tag selbst}
\begin{checkliste}
  \item Batterietest unmittelbar vor der Stunde, nicht am Vorabend.
  \item Das Brett auf einen Tisch stellen, um den die Gruppe stehen
    kann. Bei großen Gruppen eine Dokumentenkamera auf das Brett richten
    und die Exploration reihum an drei bis vier Personen delegieren.
  \item Tafel in drei Felder teilen: \enquote{Unser Modell},
    \enquote{Bestätigt durch}, \enquote{Was würde es widerlegen?}. Das
    dritte Feld bleibt in der ersten Runde absichtlich leer.
  \item Abschluss-Sprechtext (Grenzarbeit, sinngemäß): \enquote{Ich
    verrate Ihnen nicht, wie es funktioniert. Aber ich verrate Ihnen,
    wo Sie suchen müssten: Die Lösung liegt im Gerät. Es ist einfacher
    gebaut, als Sie denken, und komplizierter, als Ihr Modell erlaubt.}
    Keine weitere Andeutung; Nachfragen nach der Stunde freundlich
    abweisen.
\end{checkliste}

\textbf{Risiken und Plan B.} Fällt eine Lampe aus, wird das nicht
überspielt: Ein Defekt ist eine echte Beobachtung, die jedes Modell
erklären muss, und die Gruppe darf sie protokollieren. Ist das Gerät
ganz unbrauchbar, beginnt die Einheit mit dem 2-4-6-Test als
Staunmoment; die Zeit der Exploration (15~Min.) geht dann in einen
ausführlicheren 2-4-6-Test und die Klassifikationsübung. Kommt jemand der
Methode erkennbar nahe, weder bestätigen noch dementieren, sondern die
Vermutung wie jede andere an die Tafel schreiben und fragen, wie man sie
prüfen würde.
\end{vorbereitung}

\begin{ablauf}
\begin{description}
\item[0--5 · Blitzlicht.] Aufgreifen von LE~2: \enquote{Ihre
  Forschungsfrage ist eine Frage. Heute wird daraus eine Behauptung, die
  sich irren kann.}
\item[5--20 · Staunmoment.] Switchboard: Die Gruppe exploriert das Brett,
  protokolliert Eingaben und Beobachtungen und formuliert ein Modell.
  Dann die Behauptung der Lehrperson zur Farbe der Hülsen, Umstecken,
  neue Exploration; zwei bis drei Modellwechsel. Schlüsselfrage:
  \enquote{Welche Eingabe hätte Ihr Modell widerlegen können?}
\item[20--50 · Theorie-Input.] Induktion, Deduktion, Abduktion; Popper und
  die Falsifikation; Null- und Alternativhypothese; Horwichs graduelle
  Bestätigung und das Widerlegungsrisiko; Grenzen (Duhem-Quine).
\item[50--60 · Brücke zum Bias.] Wason-2-4-6 als Kurzübung: Die Gruppe
  bestätigt, statt zu prüfen.
\item[60--80 · Anwendung.] Klassifikationsübung mit fünf Aussagen
  (Übungsblatt Teil~A, Aufgabe~1); dann Übungsblatt Teil~B austeilen;
  jede Person formuliert für die eigene Frage H\textsubscript{0} und
  H\textsubscript{1} samt Widerlegungsbedingung (Aufgabe~3, Beginn in der
  Stunde).
\item[80--90 · Sicherung.] Kriterienraster \enquote{gute Hypothese};
  Ausblick auf Testbarkeit und Forschungsdesign (LE~6).
\end{description}
\end{ablauf}

%==============================================================================
\section{Staunmoment: Das Schaltbrett}
%==============================================================================
\begin{staunbox}{Switchboard}
Vier Schalter mit farbigen Hülsen, vier Lampen in denselben Farben.
Die Leiterbahnen sind sichtbar.

\emph{Finden Sie heraus, wie das Brett funktioniert. Protokollieren Sie
jede Eingabe und jede Beobachtung. Formulieren Sie ein Modell.}

\medskip
Danach: \emph{Ich behaupte, die Leiterbahnen sind Dekoration. Entscheidend
ist allein die Farbe der Hülse.}
\end{staunbox}
\os{\newpage}

\begin{description}
\item[{\rot\bf Runde 1: Ihr Modell}]~\\[-\lblskip]
\bi
  \ite Welches Modell haben Sie gebildet?
  \ite Wie viele Versuche haben es bestätigt?
  \ite Hat irgendjemand eine Eingabe gewählt, von der er \emph{erwartet}
    hat, dass sie das Modell bricht?
  \ite Tafel: \emph{Modell} -- \emph{Bestätigt durch} -- \emph{Was würde es
    widerlegen?}
\ei
\os{\newpage}

\item[{\rot\bf Runde 2 und 3: neue Modelle}]~\\[-\lblskip]
\bi
  \ite nach dem Umstecken ist Ihr altes Modell widerlegt -- durch
    \emph{eine} Beobachtung
  \ite das neue Modell (\enquote{Farbe entscheidet}) wird ebenso oft
    bestätigt wie das alte
  \ite und hält ebenfalls nicht
  \ite Was wissen Sie jetzt \emph{sicher} über dieses Brett?
\ei
\os{\newpage}
\end{description}

\nsl{Der Ablauf folgt der Beschreibung, die Markus Pfeil in seinem
Kommentar zu Exposure und Grenzarbeit gibt (\buch{Anhang~E}): Die
Studierenden explorieren das Brett, bilden eine Hypothese, testen sie und
ziehen ein Fazit. Danach erklärt die Lehrperson, die Verdrahtung habe mit
dem Schalten nichts zu tun, entscheidend seien die Farben von Hülse und
Lampe, und steckt die Hülsen um: Nun schaltet jeder Schalter nicht mehr
die Lampe, mit der er sichtbar verbunden ist, sondern die gleichfarbige.
Dieser Wechsel wird einige Male wiederholt. Was bleibt, ist die Erfahrung,
dass schon ein einfaches Schaltbrett zu kompliziert ist, um seine
Funktionsweise aus Bestätigungen zu erraten, und dass sich nichts davon
beweisen lässt.

Wichtig ist der Abschluss. Bleibt der Effekt ganz unkommentiert, schreiben
Studierende ihn erfahrungsgemäß der Person zu, nicht dem Gerät, und
verlieren damit die eigentliche Lehre: dass Systeme meist komplexer sind,
als wir ihnen zugestehen. Der Hinweis, dass die Lösung im Gerät liegt,
verschiebt die Wahrnehmung vom Mysterium zum Rätsel, ohne die Methode
preiszugeben.}

\begin{description}
\item[{\rot\bf Wo liegt die Lösung?}]~\\[-\lblskip]
\bi
  \ite ich verrate die Methode nicht
  \ite aber die Richtung: \textbf{die Lösung liegt im Gerät}
  \ite Ihr Modell war zu einfach für ein Brett mit vier Schaltern
  \ite Wie sicher sind Sie bei Systemen, die komplizierter sind als dieses?
\ei
\os{\newpage}
\end{description}

%==============================================================================
\section{Drei Arten zu schließen}
%==============================================================================
\nsl{Beim Explorieren des Bretts haben die Gruppen drei verschiedene
Schlussformen benutzt, meist ohne es zu bemerken. Aus einzelnen
Beobachtungen (\enquote{Schalter Rot schaltet Lampe Rot}, mehrfach) wurde
eine allgemeine Regel: das ist \textbf{Induktion}. Aus der Regel wurde
vorhergesagt, was beim nächsten Versuch passieren muss: das ist
\textbf{Deduktion}. Und als etwas Unerwartetes geschah, wurde eine
Erklärung gesucht, die das Unerwartete verständlich machen würde
(\enquote{Vielleicht gibt es eine zweite Leiterbahn auf der Rückseite}):
das ist \textbf{Abduktion}, der Schluss auf die beste Erklärung, wie ihn
Charles S.\ Peirce beschrieben hat. Wissenschaftliches Arbeiten braucht
alle drei. Abduktion erzeugt Hypothesen, Deduktion leitet aus ihnen
prüfbare Vorhersagen ab, und die Prüfung entscheidet, ob die Hypothese
bleibt. Logisch zwingend ist nur die Deduktion; Induktion und Abduktion
können bei wahren Prämissen zu falschen Schlüssen führen.}
\begin{description}
\item[{\rot\bf Induktion -- Deduktion -- Abduktion}]~\\[-\lblskip]
\begin{center}
\renewcommand{\arraystretch}{1.25}
\begin{tabular}{p{0.16\textwidth}p{0.29\textwidth}p{0.4\textwidth}}
\toprule
& \textbf{Schema} & \textbf{am Schaltbrett} \\
\midrule
\textbf{Induktion} & viele Fälle $\to$ Regel & \enquote{Fünfmal schaltete Rot die rote Lampe, also immer.} \\
\textbf{Deduktion} & Regel $\to$ Vorhersage & \enquote{Gelten die Bahnen, muss Schalter~2 Lampe~2 schalten.} \\
\textbf{Abduktion} & Überraschung $\to$ beste Erklärung & \enquote{Es gibt wohl eine verborgene Verbindung.} \\
\bottomrule
\end{tabular}
\end{center}
\bi
  \ite logisch zwingend ist nur die Deduktion
\ei
\os{\newpage}

\item[{\rot\bf Der Forschungszyklus}]~\\[-\lblskip]
\begin{center}
\resizebox{0.8\textwidth}{!}{%
\begin{tikzpicture}[font=\sffamily,>=Stealth,
  st/.style={draw=rwuviolet,thick,fill=rwuviolet!8,rounded corners=3pt,
    minimum width=34mm,minimum height=12mm,align=center}]
  \node[st] (b) at (0,0) {Beobachtung\\Überraschung};
  \node[st] (h) at (7,0) {Hypothese};
  \node[st] (v) at (7,-3) {Vorhersage};
  \node[st,draw=rwucyan,fill=rwucyan40!40] (t) at (0,-3) {Test};
  \draw[->,very thick] (b) -- node[above]{Abduktion} (h);
  \draw[->,very thick] (h) -- node[right]{Deduktion} (v);
  \draw[->,very thick] (v) -- node[below]{Prüfung} (t);
  \draw[->,very thick,rwugray] (t) -- node[left,align=right,xshift=-2mm]{bewährt oder\\verworfen} (b);
\end{tikzpicture}}
\end{center}
\bi
  \ite Hypothesen \emph{entstehen} durch Abduktion, sie \emph{bewähren}
    sich in Tests
\ei
\os{\newpage}

\item[{\rot\bf Das Induktionsproblem}]~\\[-\lblskip]
\bi
  \ite kein Stapel von Bestätigungen macht eine allgemeine Aussage
    \emph{gewiss}
  \ite tausend weiße Schwäne beweisen nicht, dass alle Schwäne weiß sind
  \ite ein einziger schwarzer Schwan widerlegt es
  \ite am Schaltbrett: viele Bestätigungen, dann ein Umstecken -- und
    das Modell ist weg
\ei
\os{\newpage}
\end{description}

%==============================================================================
\section{Popper: Falsifikation statt Verifikation}
%==============================================================================
\nsl{Karl Popper zog aus dem Induktionsproblem in der \emph{Logik der
Forschung} eine radikale Konsequenz \parencite{popper1934logik}: Wenn sich
allgemeine Aussagen nicht verifizieren lassen, dann besteht ihre Prüfung
im ernsthaften Versuch, sie zu widerlegen. Eine Aussage ist
wissenschaftlich, wenn sie prinzipiell an der Erfahrung scheitern kann.
Sie gilt als \emph{bewährt}, solange sie strengen Widerlegungsversuchen
standhält, und bleibt dabei vorläufig.

Für das Buch ist entscheidend, dass Popper damit weniger eine Eigenschaft
von Sätzen als eine Haltung von Menschen beschreibt (\buch{Kap.~5b}). Eine
Theorie wird nicht dadurch unwissenschaftlich, dass sie falsch ist; falsch
sein darf jede. Unwissenschaftlich wird sie, wenn ihre Vertreter sie gegen
jede Widerlegung abdichten. Popper nannte das \emph{Immunisierung}. Sein
Paradebeispiel war die Astrologie, deren Vertreter ihre Deutungen so vage
hielten, dass nichts mehr als Widerlegung zählen konnte
\parencite{popper1963conjectures}. Dieses Beispiel steht übrigens nicht in
der \emph{Logik der Forschung}, sondern in \emph{Conjectures and
Refutations}. Die falsche Zuschreibung ist ein verbreiteter Zitierfehler
und ein gutes Argument dafür, Primärquellen selbst zu lesen (LE~4).}
\begin{description}
\item[{\rot\bf Poppers Kriterium}]~
\end{description}
\defn{falsifizierbar}{Eine Aussage ist \textbf{falsifizierbar}, wenn
sich mindestens eine mögliche Beobachtung angeben lässt, die ihr
widerspräche. Sie gilt als \textbf{bewährt}, solange ernsthafte
Widerlegungsversuche gescheitert sind. Bewiesen ist sie damit nicht.}
\bi
  \ite falsifizierbar heißt nicht falsch, sondern \emph{prüfbar}
  \ite Gegenbegriff: Immunisierung -- die Aussage wird gegen jede
    Widerlegung abgedichtet
\ei
\os{\newpage}
\begin{description}

\item[{\rot\bf Die Asymmetrie}]~\\[-\lblskip]
\begin{center}
\resizebox{0.8\textwidth}{!}{%
\begin{tikzpicture}[font=\sffamily,>=Stealth,
  box/.style={draw=rwuviolet,thick,fill=rwuviolet!8,rounded corners=2pt,
    minimum width=40mm,minimum height=11mm,align=center}]
  \node[box] (h) at (0,0) {Hypothese $H$\\\enquote{alle Schwäne sind weiß}};
  \node[box,draw=rwucyan,fill=rwucyan40!40] (j) at (-4.5,-2.8) {1000 weiße Schwäne};
  \node[box,draw=rwucyan,fill=rwucyan40!40] (n) at (4.5,-2.8) {1 schwarzer Schwan};
  \draw[->,thick] (j) -- (h);
  \draw[->,thick,rwuviolet] (n) -- (h);
  \node[below=2mm of j,align=center] {$H$ wird wahrscheinlicher};
  \node[below=2mm of n,align=center,text=rwuviolet] {$H$ ist widerlegt};
\end{tikzpicture}}
\end{center}
\bi
  \ite Bestätigung ist graduell, Widerlegung kann schlagend sein
\ei
\os{\newpage}

\item[{\rot\bf Riskante Vorhersagen}]~\\[-\lblskip]
\bi
  \ite eine Vorhersage, die auch bei falscher Theorie einträfe, prüft
    nichts
  \ite je genauer und unwahrscheinlicher die Vorhersage, desto mehr ist
    ihr Eintreffen wert
  \ite Beispiel: die Vorhersage der Position Neptuns 1846 -- für einen
    Körper, den noch niemand gesehen hatte
  \ite Gegenbeispiel: ein Horoskop, in dem sich jede Leserin wiederfindet
\ei
\os{\newpage}

\item[{\rot\bf Immunisierung erkennen}]~\\[-\lblskip]
\bi
  \ite \enquote{Es wirkt, aber nicht unter Beobachtung.}
  \ite \enquote{Wenn es nicht klappt, war der Glaube zu schwach.}
  \ite Begriffe, die nachträglich so gedehnt werden, dass alles passt
  \ite Test: \emph{Welches Ergebnis würden Sie als Widerlegung
    akzeptieren?} Gibt es keins, ist die Aussage nicht prüfbar
\ei
\os{\newpage}
\end{description}

%==============================================================================
\section{Von der Vermutung zur Hypothese}
%==============================================================================
\nsl{Eine Vermutung ist ein Einfall. Eine Hypothese ist eine Vermutung,
die so präzise gefasst ist, dass man sie prüfen kann. Drei Eigenschaften
machen den Unterschied: \emph{Präzision} (Was genau wird behauptet, über
wen, unter welchen Bedingungen?), \emph{Testbarkeit} (Welche Beobachtung
würde widersprechen?) und \emph{Reichweite} (Für welche Fälle soll die
Aussage gelten?). Präzision entsteht durch \emph{Operationalisierung}: Die
Begriffe der Hypothese werden mit beobachtbaren Messvorschriften
verbunden. Das Konzept ist aus LE~2 bekannt; hier wird es zur Bedingung
der Falsifizierbarkeit. \enquote{Ehrliche Menschen schlafen besser} ist in
dieser Form kaum prüfbar, weil weder \enquote{ehrlich} noch \enquote{besser
schlafen} festgelegt ist. Erst mit einer Messvorschrift für beides,
etwa einem etablierten Fragebogen zur Ehrlichkeit und der mit einem
Aktivitätstracker gemessenen Schlafdauer über zwei Wochen, entsteht eine
Aussage, die scheitern kann.}
\begin{description}
\item[{\rot\bf Vermutung oder Hypothese?}]~\\[-\lblskip]
\begin{center}
\renewcommand{\arraystretch}{1.25}
\begin{tabular}{p{0.18\textwidth}p{0.31\textwidth}p{0.36\textwidth}}
\toprule
& \textbf{Vermutung} & \textbf{Hypothese} \\
\midrule
Präzision & \enquote{Ehrliche schlafen besser.} & Wer ist ehrlich? Was heißt besser? -- festgelegt \\
Testbarkeit & offen & benannte Widerlegungsbedingung \\
Reichweite & unklar & Population, Bedingungen, Zeitraum \\
\bottomrule
\end{tabular}
\end{center}
\bi
  \ite Operationalisierung macht aus Begriffen Messvorschriften (LE~2)
\ei
\os{\newpage}

\item[{\rot\bf Arten von Hypothesen}]~\\[-\lblskip]
\bi
  \ite \textbf{Unterschiedshypothese:} Gruppe A unterscheidet sich von
    Gruppe B in Merkmal Y
  \ite \textbf{Zusammenhangshypothese:} X und Y hängen zusammen
  \ite \textbf{Veränderungshypothese:} Y ändert sich zwischen
    Messzeitpunkt 1 und 2
  \ite \textbf{gerichtet} (\enquote{A $>$ B}) oder \textbf{ungerichtet}
    (\enquote{A $\neq$ B})
\ei
\os{\newpage}
\end{description}

%==============================================================================
\section{Null- und Alternativhypothese}
%==============================================================================
\nsl{In der empirischen Forschung, vor allem in der quantitativen, wird
die eigene Vermutung als \emph{Alternativhypothese} H\textsubscript{1}
formuliert. Ihr gegenüber steht die \emph{Nullhypothese}
H\textsubscript{0}, die behauptet, dass es den vermuteten Unterschied oder
Zusammenhang nicht gibt. Geprüft wird formal die Nullhypothese: Die Daten
werden daraufhin befragt, wie gut sie mit ihr vereinbar sind. Das wirkt
umständlich, hat aber einen guten Grund. Die Nullhypothese ist präzise
(\enquote{Unterschied gleich null}), die Alternativhypothese meist nicht
(\enquote{irgendein Unterschied}), und präzise Aussagen lassen sich gezielt
unter Druck setzen. Die Rechentechnik folgt in LE~8. Hier geht es darum,
beide Hypothesen sauber zu formulieren und vorab festzulegen, welches
Ergebnis zur Verwerfung der eigenen Hypothese führen würde
\parencite[vgl.][]{doering2016forschungsmethoden}.}
\begin{description}
\item[{\rot\bf H\textsubscript{0} und H\textsubscript{1}}]~\\[-\lblskip]
\bi
  \ite \textbf{H\textsubscript{1}} (Alternativhypothese): Ihre Vermutung --
    es gibt einen Unterschied oder Zusammenhang
  \ite \textbf{H\textsubscript{0}} (Nullhypothese): es gibt ihn nicht
  \ite beide schließen sich aus und decken zusammen alle Fälle ab
  \ite geprüft wird H\textsubscript{0}: Sind die Daten mit \enquote{kein
    Effekt} gut vereinbar?
  \ite Rechnen kommt in LE~8 -- heute: sauber formulieren
\ei
\os{\newpage}

\item[{\rot\bf Beispiel}]~\\[-\lblskip]
\bi
  \ite Frage: Lernen Studierende mit Staunmoment mehr?
  \ite H\textsubscript{1}: Wer vor einem Lehrvideo einen Zaubertrick
    sieht, löst eine Transferaufgabe häufiger richtig als wer ihn nicht
    sieht
  \ite H\textsubscript{0}: Die Lösungsquoten unterscheiden sich nicht
  \ite Widerlegungsbedingung (vorab festgelegt): gleiche oder niedrigere
    Quote in der Trickgruppe
\ei
\os{\newpage}
\end{description}

\nsl{Das Beispiel ist nicht erfunden: Es ist die Leitfrage der
Becher-Monty-Hall-Studien aus dem Buch (\buch{Kap.~12d und 12f}), die in
LE~6 als Fallbeispiel für Forschungsdesign wiederkehren. Dort zeigt sich
auch, wie wichtig eine vorab festgelegte Widerlegungsbedingung ist: Ein
Zwischenstand der Feldstudie wies in die erhoffte Richtung und verschwand
mit wachsender Stichprobe.}

%==============================================================================
\section{Horwich: Bestätigung hat Grade}
%==============================================================================
\nsl{Poppers Asymmetrie lässt sich wahrscheinlichkeitstheoretisch genauer
fassen. Paul Horwich hat in \emph{Probability and Evidence} gezeigt, dass
viele Rätsel der Wissenschaftstheorie verschwinden, sobald man Überzeugung
als graduell behandelt \parencite{horwich1982probability}. Folgt eine
Beobachtung $E$ logisch aus einer Hypothese $H$ und tritt sie nicht ein,
dann ist $H$ widerlegt: $P(H\mid\neg E)=0$. Tritt $E$ dagegen ein, gilt
nur $P(H\mid E)>P(H)$. Die Hypothese wird wahrscheinlicher, bewiesen ist
sie nicht. Wie viel wahrscheinlicher, hängt davon ab, wie unwahrscheinlich
$E$ ohnehin war: Impliziert $H$ die Beobachtung $E$, ist der
Bestätigungsfaktor $P(H\mid E)/P(H) = 1/P(E)$. Eine Vorhersage, die
sowieso jeder erwartet hätte, bringt also fast nichts.

Horwichs Auflösung des Rabenparadoxons führt das an einem Bild vor
(\buch{Kap.~5b, Horwichs Theorie der Evidenz}): Wer einen Raben prüft
und ihn schwarz findet, hat die Hypothese \enquote{alle Raben sind schwarz}
einem Risiko ausgesetzt. Wer ein schwarzes Objekt prüft und feststellt,
dass es ein Rabe ist, hat nichts riskiert. Nur Beobachtungen, die ein
Widerlegungsrisiko eingehen, haben vollen Beweiswert.}
\begin{description}
\item[{\rot\bf Bestätigung ist graduell}]~\\[-\lblskip]
\bi
  \ite tritt die vorhergesagte Beobachtung \emph{nicht} ein:
    $P(H\mid\neg E)=0$ -- widerlegt
  \ite tritt sie ein: nur $P(H\mid E) > P(H)$ -- wahrscheinlicher, nicht
    bewiesen
  \ite Stärke der Bestätigung: $P(H\mid E)\,/\,P(H) = 1/P(E)$
  \ite \textbf{je unwahrscheinlicher die Vorhersage, desto stärker die
    Bestätigung}
\ei
\os{\newpage}

\item[{\rot\bf Das Widerlegungsrisiko}]~\\[-\lblskip]
\bi
  \ite einen \emph{Raben} prüfen und schwarz finden: echter Test
  \ite etwas \emph{Schwarzes} prüfen und einen Raben finden: kein Test
  \ite am Schaltbrett: \enquote{Rot schaltet Rot} zum zehnten Mal prüfen
    ist der schwarze Gegenstand
  \ite Regel: \emph{Wählen Sie den Versuch, dessen Ausgang Sie am
    wenigsten sicher vorhersagen können.}
\ei
\os{\newpage}
\end{description}

\begin{merkbox}
Bestätigungen sammeln ist billig. Eine Beobachtung ist nur so viel wert
wie das Risiko, das die Hypothese mit ihr eingegangen ist. Fragen Sie
deshalb nicht \enquote{Was spricht für meine Hypothese?}, sondern
\enquote{Was würde mich umstimmen?}
\end{merkbox}
\os{\newpage}

%==============================================================================
\section{Grenzen der Falsifikation: Duhem-Quine}
%==============================================================================
\nsl{Auch die Widerlegung ist weniger eindeutig, als Poppers Schwan
suggeriert. Pierre Duhem hat für die Physik gezeigt, dass eine Hypothese
nie allein geprüft wird, sondern immer zusammen mit Hilfsannahmen über
Messgeräte, Versuchsaufbau und Randbedingungen \parencite{duhem1906theorie}.
Widerspricht ein Experiment der Vorhersage, sagt die Logik nur, dass
\emph{irgendein} Glied der Kette falsch ist, nicht welches. Willard Van
Orman Quine hat den Gedanken auf das gesamte Wissenssystem ausgeweitet
\parencite{quine1951dogmas}. Am Schaltbrett ist das unmittelbar zu sehen:
Leuchtete eine Lampe nicht wie vorhergesagt, konnte man das Modell
verwerfen oder eine Hilfsannahme (\enquote{Lampe defekt}, \enquote{Batterie
schwach}, \enquote{falsch geschaltet}). Daraus folgt nicht, dass
Falsifikation wertlos ist. Es folgt, dass Hilfsannahmen offengelegt und, wo
möglich, unabhängig geprüft werden müssen. Ad-hoc-Rettungen, die nur dazu
dienen, eine Hypothese gegen ein unbequemes Ergebnis zu schützen, sind
genau die Immunisierung, vor der Popper warnt. Imre Lakatos hat daraus ein
Kriterium gemacht: Ein Forschungsprogramm schreitet fort, wenn seine
Anpassungen neue, prüfbare Vorhersagen liefern, und degeneriert, wenn sie
nur Löcher stopfen \parencite{lakatos1970falsification}.}
\begin{description}
\item[{\rot\bf Duhem-Quine-These}]~\\[-\lblskip]
\bi
  \ite geprüft wird nie $H$ allein, sondern $H$ \emph{plus} Hilfsannahmen
    $A_1, A_2, \ldots$ (Messgerät, Aufbau, Randbedingungen)
  \ite widerspricht das Ergebnis, ist \emph{irgendetwas} falsch --
    aber was?
  \ite Schaltbrett: Modell falsch -- oder Lampe defekt? Batterie
    schwach?
  \ite Konsequenz: Hilfsannahmen offenlegen und eigens prüfen
\ei
\os{\newpage}

\item[{\rot\bf Rettung oder Ausrede?}]~\\[-\lblskip]
\bi
  \ite legitim: eine Hilfsannahme, die man \emph{unabhängig} prüfen kann
    (Lampe tauschen, Versuch wiederholen)
  \ite Ausrede: eine Zusatzannahme, die nur dazu dient, das Ergebnis
    wegzuerklären
  \ite Lakatos: fortschreitende Forschungsprogramme sagen durch ihre
    Anpassungen \emph{Neues} voraus, degenerierende nicht
\ei
\os{\newpage}
\end{description}

%==============================================================================
\section{Brücke zum Bias: der 2-4-6-Test}
%==============================================================================
\nsl{Dass die Gruppe am Schaltbrett fast nur bestätigende Eingaben
gewählt hat, ist kein Zufall. Peter Wason hat diese Neigung 1960 mit einer
einfachen Zahlenaufgabe sichtbar gemacht \parencite{wason1960failure}: Die
Versuchspersonen sollten die Regel hinter dem Tripel \enquote{2--4--6}
finden, indem sie eigene Tripel vorschlugen. Die Regel war weit gefasst
(\enquote{drei aufsteigende Zahlen}); die meisten bildeten eine engere
Hypothese und testeten nur Tripel, die dazu passten. Nur sechs von 29
Personen fanden die Regel ohne falschen Zwischenschluss. Jonathan Baron
nennt die Neigung, eine Hypothese ausschließlich direkt zu prüfen,
\emph{Congruence Bias} \parencite{baron1988diagnostic}. Fair bleibt der
Hinweis von \textcite{klaymanha1987}: Die \emph{positive Teststrategie} ist
in vielen Umgebungen eine gute Heuristik. Sie versagt dort, wo die wahre
Regel weiter ist als die eigene Hypothese, und das kommt in der Forschung
häufig vor. Der Ablauf unten folgt der Unterrichtsfassung des Buchs
(\buch{Kap.~10, Der 2-4-6-Test}), auf zehn Minuten gekürzt.}
\begin{description}
\item[{\rot\bf Die Regel finden}]~\\[-\lblskip]
\begin{center}
{\Huge\bfseries\color{rwuviolet} 2 \quad 4 \quad 6}
\end{center}
\bi
  \ite diese drei Zahlen folgen einer Regel, die ich mir ausgedacht habe
  \ite schlagen Sie eigene Tripel vor -- ich sage nur \emph{Ja} oder
    \emph{Nein}
  \ite nennen Sie die Regel erst, wenn Sie ganz sicher sind
\ei
\os{\newpage}
\end{description}

\begin{lehrhinweis}[Ablauf 2-4-6-Test (10 Min.)]
\begin{description}
  \item[1 Min.] \enquote{2--4--6} an die Tafel, Auftrag wie auf der Folie.
  \item[4 Min.] Vorschläge aus dem Plenum, jeweils mit \emph{Ja} oder
    \emph{Nein} in zwei Spalten an der Tafel protokollieren. Typisch sind
    8--10--12, 20--22--24, 100--102--104: alles \emph{Ja}.
  \item[1 Min.] Jede Person notiert still ihre vermutete Regel auf einem
    Zettel.
  \item[2 Min.] Frage: \enquote{Hat jemand ein Tripel vorgeschlagen, bei
    dem er ein \emph{Nein} erwartet hat?} Meist niemand. Dann gezielt
    widersprechende Tripel testen lassen: 1--2--3 (\emph{Ja}), 3--2--1
    (\emph{Nein}), 2--4--5 (\emph{Ja}).
  \item[2 Min.] Auflösung: \enquote{drei aufsteigende Zahlen}. Rückbezug
    auf das Schaltbrett und auf Horwichs Raben: Die bestätigenden Tripel
    waren schwarze Gegenstände.
\end{description}
\end{lehrhinweis}

\begin{description}
\item[{\rot\bf Was ist passiert?}]~\\[-\lblskip]
\bi
  \ite die Regel war: \textbf{drei aufsteigende Zahlen}
  \ite fast alle Vorschläge passten zur eigenen, engeren Hypothese -- und
    bekamen ein \emph{Ja}
  \ite Wason 1960: nur 6 von 29 Personen fanden die Regel ohne falschen
    Zwischenschluss
  \ite \textbf{Congruence Bias}: man testet nur die eigene Hypothese, nie
    die Alternativen
  \ite Gegenmittel (Baron): Alternative benennen, dann den Test wählen,
    der zwischen beiden \emph{unterscheidet}
\ei
\os{\newpage}
\end{description}

%==============================================================================
\section{Übung in der Stunde}
%==============================================================================
\begin{description}
\item[{\rot\bf Falsifizierbar oder nicht? (Kleingruppen, 10 Min.)}]~\\[-\lblskip]
\bi
  \ite Übungsblatt Teil~A, Aufgabe~1: fünf Aussagen
  \ite Welche sind falsifizierbar -- so, wie sie dastehen?
  \ite für jede prüfbare Aussage: eine konkrete Widerlegungsbedingung
  \ite für jede vage Aussage: eine präzisierte, prüfbare Fassung
\ei
\os{\newpage}

\item[{\rot\bf Ihre eigene Hypothese (Einzelarbeit, 10 Min.)}]~\\[-\lblskip]
\bi
  \ite nehmen Sie Ihre Forschungsfrage aus LE~2
  \ite formulieren Sie H\textsubscript{1} und H\textsubscript{0}
  \ite Welche Begriffe müssen Sie operationalisieren?
  \ite Bei welchem Ergebnis würden Sie H\textsubscript{1} verwerfen?
  \ite Tausch mit der Nachbarin: Würde \emph{sie} Ihre
    Widerlegungsbedingung akzeptieren?
\ei
\os{\newpage}
\end{description}

\begin{lehrhinweis}
Bei der Klassifikationsübung lohnt es, auf zwei Fehler zu achten.
Erstens wird \enquote{falsifizierbar} oft mit \enquote{falsch} oder
\enquote{unbewiesen} verwechselt; hier hilft die Frage, welche Beobachtung
der Aussage widersprechen würde. Zweitens halten viele Gruppen vage
Aussagen pauschal für nicht falsifizierbar. Die bessere Antwort ist meist:
\enquote{In dieser Form nicht prüfbar, aber präzisierbar.} Bei der eigenen
Hypothese zeigt der Partnertausch schnell, ob die Widerlegungsbedingung
beobachtbar ist. Formulierungen wie \enquote{wenn es nicht stimmt} oder
\enquote{wenn kein Zusammenhang besteht} ohne Messung und Kriterium reichen
nicht. Der vollständige Erwartungshorizont steht am Ende von Teil~B des
Übungsblatts.
\end{lehrhinweis}

%==============================================================================
\section{Sicherung: Was ist eine gute Hypothese?}
%==============================================================================
\begin{description}
\item[{\rot\bf Kriterienraster}]~\\[-\lblskip]
\begin{center}
\renewcommand{\arraystretch}{1.2}
\begin{tabular}{p{0.24\textwidth}p{0.64\textwidth}}
\toprule
\textbf{Kriterium} & \textbf{Prüffrage} \\
\midrule
präzise & Sind alle Begriffe operationalisiert? \\
falsifizierbar & Welches Ergebnis würde ihr widersprechen? \\
riskant & Wäre das Ergebnis ohne die Hypothese unwahrscheinlich? \\
begrenzt & Für wen, wann, unter welchen Bedingungen gilt sie? \\
begründet & Woraus folgt sie (Theorie, Literatur, Beobachtung)? \\
offen & Sind die Hilfsannahmen benannt? \\
\bottomrule
\end{tabular}
\end{center}
\os{\newpage}
\end{description}

\begin{merkbox}
Eine Hypothese ist gut, wenn sie scheitern \emph{kann} und Sie vorher
sagen können, woran. Bewährt ist sie, wenn sie solche Versuche
überstanden hat. Bewiesen ist sie nie.
\end{merkbox}
\os{\newpage}

%==============================================================================
\section{Lernziele und Aufgaben}
%==============================================================================
\begin{description}
\item[{\rot\bf Nach dieser Einheit können Sie}]~\\[-\lblskip]
\bi
  \ite Induktion, Deduktion und Abduktion unterscheiden
  \ite Hypothesen nach dem Falsifikationskriterium bewerten
  \ite empirisch prüfbare von nicht prüfbaren Aussagen unterscheiden
  \ite eigene Hypothesen operationalisierbar formulieren, mit
    H\textsubscript{0}, H\textsubscript{1} und Widerlegungsbedingung
  \ite erklären, warum Bestätigung graduell ist und Widerlegung an
    Hilfsannahmen hängt
\ei
\os{\newpage}

\item[{\rot\bf Bis zur nächsten Woche}]~\\[-\lblskip]
\bi
  \ite \textbf{Übungsblatt}, Aufgabe~1 (Teil~A) und 2 (Teil~B) fertigstellen
  \ite \textbf{schriftlich} (Aufgabe~3): H\textsubscript{0} und
    H\textsubscript{1} zu Ihrer Forschungsfrage, mit den Bedingungen, unter
    denen Sie Ihre Hypothese verwerfen würden
  \ite \textbf{Tagebuch des Staunens:} ein Erlebnis dieser Woche, bei dem
    Sie etwas nur bestätigt statt geprüft haben
  \ite Vorbereitung LE~6: Wie würden Sie Ihre Hypothese testen?
\ei
\os{\newpage}
\end{description}

\handoutinput{handouts/LE05}
